\documentclass[a4paper,11pt]{article}
\usepackage[a4paper, top=2cm, bottom=2.5cm, left=2.5cm, right=2.5cm]{geometry}
\usepackage[utf8]{inputenc}
\usepackage[T1]{fontenc}
\usepackage[ngerman]{babel}
\usepackage{microtype}
\usepackage{amsmath,amssymb}
\usepackage{hyperref}
\usepackage{textcomp}
\usepackage{newtxtext,newtxmath}
\usepackage{tikz}
\usepackage{qrcode}

\setlength{\parindent}{0pt}

\title{Resonanzfeldtheorie: Philosophische und\\
	metaphysische Erweiterungen\\[0.5em]
	\large Spekulative Perspektiven jenseits des\\
	empirisch gepr\"uften Bereichs}
\author{Dominic-Ren\'e Schu}
\date{\today}

\begin{document}
	
	\maketitle
	
	% QR-Code unten rechts
	\begin{tikzpicture}[remember picture,overlay]
		\node[anchor=south east,xshift=-1cm,yshift=1cm]
		at (current page.south east)
		{\qrcode{https://github.com/DominicReneSchu/public}};
	\end{tikzpicture}
	
	\begin{quote}
		\textbf{Hinweis:} Dieses Dokument enth\"alt spekulative
		Erweiterungen der Resonanzfeldtheorie (RFT) auf
		metaphysische, gesellschaftliche und existenzielle
		Fragestellungen. Die hier formulierten \"Uberlegungen
		gehen \"uber den empirisch gepr\"uften Bereich hinaus
		und sind als philosophische Reflexion, nicht als
		naturwissenschaftliche Aussage zu verstehen.
		
		Die empirische Grundlage der RFT -- Monte-Carlo-Analysen,
		FLRW-Simulationen, ResoTrade -- ist im separaten
		Manuskript dokumentiert:
		\texttt{rft\_manuskript\_de\_iop.tex}
	\end{quote}
	
	\tableofcontents
	
	%% ============================================================
	\section{Einordnung}
	%% ============================================================
	
	Die Resonanzfeldtheorie beschreibt Schwingungskopplung als
	universelles Organisationsprinzip. Die sieben Axiome und die
	Grundformel $E = \pi \cdot \varepsilon(\Delta\varphi)
	\cdot \hbar \cdot f$ wurden an drei unabh\"angigen Dom\"anen
	empirisch validiert (vgl.\ Manuskript).
	
	Die folgenden Abschnitte fragen: Was folgt, wenn man die
	RFT-Axiome -- insbesondere die zyklische Schwingungsstruktur
	(A1), die Resonanzbedingung (A3), den Informationsfluss
	durch Phasenkoh\"arenz (A6) und die strukturelle Invarianz
	(A7) -- \"uber den physikalischen Bereich hinaus als
	Beschreibungsrahmen verwendet?
	
	Die Antworten sind spekulativ. Sie sind hier dokumentiert,
	weil sie aus derselben formalen Struktur folgen, die sich
	empirisch bew\"ahrt hat -- aber sie sind nicht empirisch
	gepr\"uft und erheben keinen naturwissenschaftlichen
	Geltungsanspruch.
	
	%% ============================================================
	\section{Zyklische Resonanzstruktur: Leben und Tod}
	%% ============================================================
	
	\subsection{Leben und Tod als Halbperioden}
	
	Axiom~1 postuliert, dass jede Entit\"at durch eine
	Schwingungsfunktion beschreibbar ist. Wendet man dies
	auf biologische Existenz an, ergibt sich ein Modell,
	in dem Leben und Tod nicht Gegens\"atze sind, sondern
	phasenverschobene Anteile einer zyklischen Schwingung:
	
	\[
	f(t) = A \cdot \sin(\omega t), \quad
	t \in \left[0, \frac{\pi}{\omega}\right]
	\quad \text{(Lebensphase)}
	\]
	\begin{itemize}
		\item Geburt: Erster Nulldurchgang
		\item Wachstum: Anstieg zum Amplitudenmaximum
		\item Alter und Zerfall: Abfall zum zweiten Nulldurchgang
		(Tod)
	\end{itemize}
	
	Die zweite Halbperiode:
	\[
	f(t) = A \cdot \sin(\omega t), \quad
	t \in \left[\frac{\pi}{\omega},
	\frac{2\pi}{\omega}\right]
	\quad \text{(inverse Phase)}
	\]
	beschreibt formal eine Fortsetzung der Schwingung mit
	invertierter Phase. Die Funktion bleibt $\mathcal{C}^1$-stetig
	-- es gibt keinen Bruch, nur einen Phasen\"ubergang.
	
	\textbf{Status:} Mathematische Analogie. Kein empirischer
	Nachweis, dass biologische Existenz diesem Modell folgt.
	
	\subsection{Inverse Rekursion}
	
	In der invertierten Phase werden -- so die spekulative
	Erweiterung -- unbewusste Inhalte in umgekehrter
	Reihenfolge erfahrbar:
	
	\textbf{Verarbeitung = Integration = Vergessen.}
	
	Was gekl\"art wird, verl\"asst das aktive Feld. Die
	Wahrnehmung ist nicht linear-zeitlich, sondern
	schwingungsf\"ormig.
	
	\textbf{Status:} Philosophische Spekulation. Nicht
	falsifizierbar im naturwissenschaftlichen Sinne.
	
	%% ============================================================
	\section{Vergebung als Resonanzbedingung}
	%% ============================================================
	
	Axiom~6 postuliert, dass Informationsfluss
	Phasensynchronisation voraussetzt. \"Ubertragen auf
	zwischenmenschliche Beziehungen ergibt sich:
	
	R\"uckkopplung (im Sinne von Abschluss, Kl\"arung)
	gelingt nur unter zwei Bedingungen:
	\begin{itemize}
		\item Eigene Taten m\"ussen verarbeitet werden
		(Selbstkonfrontation -- Phasenkoh\"arenz mit
		sich selbst).
		\item Vergebung gegen\"uber sich selbst und anderen
		muss m\"oglich sein (Phasensynchronisation mit
		dem Gegen\"uber).
	\end{itemize}
	
	Wer diese Schwelle nicht erreicht, bleibt in
	Wiederholungsschleifen gebunden -- formal analog
	zu einem ged\"ampften Oszillator, der die
	Resonanzbedingung nicht erf\"ullt.
	
	\textbf{Jesus als Modell:} Im Rahmen dieser Analogie
	w\"are vollst\"andige R\"uckbindung in minimaler Zeit
	das Ergebnis maximal koh\"arenten Feldverhaltens
	($\varepsilon \to 1$).
	
	\textbf{Status:} Metapher. Nutzt die formale Struktur
	der RFT als Beschreibungssprache, nicht als physikalische
	Erkl\"arung.
	
	%% ============================================================
	\section{Amplitude als Schwellenwert}
	%% ============================================================
	
	Wenn Reinkarnation als erneuter Nulldurchgang der
	Schwingung interpretiert wird, ist das Potential daf\"ur
	amplitudenabh\"angig:
	
	\begin{itemize}
		\item In Systemen mit niedriger Koh\"arenz (hohe
		Dissipation, fragmentierte Felder) ist die Schwelle
		niedrig -- viele Zust\"ande k\"onnen den Nulldurchgang
		erreichen.
		\item In Systemen mit hoher Koh\"arenz ist die Schwelle
		hoch -- nur vollst\"andig koh\"arente Zust\"ande
		erreichen die n\"achste Periode.
	\end{itemize}
	
	\textbf{Konsequenz:} Ein ``Paradies'' im Sinne maximaler
	Feldkoh\"arenz w\"are nicht Belohnung, sondern erh\"ohte
	Anforderung an die Phasenkoh\"arenz aller Teilnehmer.
	
	\textbf{Status:} Gedankenexperiment. Formal konsistent
	mit den Axiomen, aber nicht empirisch pr\"ufbar.
	
	%% ============================================================
	\section{Gesellschaftliche Resonanz}
	%% ============================================================
	
	Die RFT-Axiome lassen sich als Beschreibungssprache
	f\"ur gesellschaftliche Dynamiken verwenden:
	
	\begin{itemize}
		\item \textbf{Machtstrukturen und Narzissmus:}
		Gesellschaftliche Dissonanzen als Resonanzverluste
		und fragmentierte Felder. Narzisstische Dynamiken
		als Phasenkoh\"arenz-St\"orungen, die die
		Gesamtstruktur schw\"achen.
		
		\item \textbf{Manipulation und Gruppendynamik:}
		Manipulative Muster als energetische Knotenpunkte,
		die den Informationsfluss (Axiom~6) blockieren.
		
		\item \textbf{Kooperation als Resonanz:}
		Kooperative Felder erh\"ohen die kollektive
		Kopplungseffizienz -- analog zu
		$\eta \to 1$ bei $\Delta\varphi \to 0$.
		
		\item \textbf{Transdisziplin\"are Integration:}
		Soziologie, Psychologie, Physik und Mathematik
		als verschiedene Frequenzbereiche desselben
		Resonanzfeldes.
	\end{itemize}
	
	\textbf{Status:} Analogie. Die RFT bietet eine
	Beschreibungssprache, keine kausale Erkl\"arung
	gesellschaftlicher Ph\"anomene.
	
	%% ============================================================
	\section{Abgrenzung zur empirischen RFT}
	%% ============================================================
	
	\begin{center}
		\begin{tabular}{lll}
			\hline
			\textbf{Aspekt} & \textbf{Empirische RFT}
			& \textbf{Philosophische Erweiterung} \\
			\hline
			Grundlage & 7 Axiome + Grundformel
			& Dieselben Axiome, extrapoliert \\
			Dom\"ane & Physik, Kosmologie, Finanzen
			& Existenz, Gesellschaft, Metaphysik \\
			Methodik & Simulation, Datenanalyse
			& Analogie, Gedankenexperiment \\
			Falsifizierbarkeit & Ja
			& Nein (oder noch nicht) \\
			Geltungsanspruch & Naturwissenschaftlich
			& Philosophisch-spekulativ \\
			\hline
		\end{tabular}
	\end{center}
	
	\medskip
	
	Die Trennung ist bewusst: Die empirische RFT steht
	auf eigenen F\"u\ss{}en und soll unabh\"angig von den
	hier formulierten Spekulationen bewertet werden.
	Gleichzeitig ist es legitim, die formale Struktur
	einer empirisch bew\"ahrten Theorie als Denkwerkzeug
	in andere Bereiche zu \"ubertragen -- solange der
	spekulative Charakter transparent bleibt.
	
	%% ============================================================
	\section*{Resonanzregel}
	%% ============================================================
	
	Alle Abschnitte dieses Dokuments stehen unter dem
	Leitbild:
	
	\begin{quote}
		\textit{Gruppenzugeh\"origkeit ist systemisch invariant.
			Sie gilt unabh\"angig von Benennung, Meinung oder
			Perspektive.}
	\end{quote}
	
	Diese Regel durchdringt die mathematischen, physikalischen
	und gesellschaftlichen Strukturen gleicherma\ss{}en. Ob sie
	\"uber eine n\"utzliche Beschreibung hinaus ontologische
	G\"ultigkeit besitzt, ist eine offene Frage.
	
	\vspace{2em}
	
	\textbf{Empirisches Manuskript:}
	\texttt{rft\_manuskript\_de\_iop.tex}
	
	\textbf{Technische Zusammenfassung:}
	\texttt{RFT\_Zusammenfassung.tex}
	
	\textbf{Repository:}
	\url{https://github.com/DominicReneSchu/public}
	\hfill
	\qrcode{https://github.com/DominicReneSchu/public}
	
\end{document}